% ============================================================
\section{Rel (Relations)}
A categoria de relações é a base para modelar qualquer tipo de conectividade genérica ou base de dados relacionais.
Um objeto nesta categoria define-se por um conjunto finito e uma relação binária arbitrária (sem exigência de reflexividade, simetria ou transitividade).

\subsection{Objetos}
No ambiente MATLAB, representamos estas estruturas através de matrizes de adjacência booleanas simples.
A ausência de restrições algébricas torna a sua instanciação direta e computacionalmente leve.

\begin{lstlisting}[caption={Representação de uma Relação Arbitrária}]
% Matriz de adjacência para uma relação genérica
n = 3;
M = [0 1 0;
0 0 1; 1 0 0];
\end{lstlisting}

\subsection{Morfismos}
Os morfismos nesta categoria são funções que preservam a relação.
Um mapeamento $f$ é válido se, para quaisquer elementos $i, j$, a existência de uma relação $i \sim j$ implicar a relação $f(i) \sim f(j)$.

\begin{lstlisting}[caption={Algoritmo de Verificação de Preservação Relacional}]
% f é o vetor de mapeamento
isMorphism = true;
[I, J] = find(M_domain);
for k = 1:length(I)
    if ~M_codomain(f(I(k)), f(J(k)))
        isMorphism = false;
        break;
    end
end
\end{lstlisting}